\par\bigskip\Needspace{5\baselineskip}
\begin{center}\bfseries 5.\end{center}
\noindent\hspace*{\parindent}Einige andere Sätze übergehend wenden wir uns nun zur Theorie der quadratischen Reste, welcher \textls[100]{der vierte Abschnitt} gewidmet ist.

Die Frage nach den quadratischen Resten eines Moduls \(m\) ist identisch mit der Frage nach der Auflösbarkeit der binomischen Kongruenz zweiten Grades
\[
x^2\equiv a\ (\mathrm{mod.}\ m).
\]

Sie bietet aber eine zwiefache Seite dar, jenachdem man den Modul \(m\) oder den Rest \(a\) als gegeben betrachtet. In der erstern Annahme verursacht die Beantwortung der Frage keine erhebliche Schwierigkeit. Es ist leicht, sie auf den Fall zurückzuführen, wo der Modul eine Primzahlpotenz ist, und, da sie für den Modul \(2^n\) einfach erledigt und für den Modul \(p^n\), wo \(p\) eine ungerade Primzahl bedeutet, aus dem Falle des einfachen Moduls \(p\) entschieden werden kann, so handelt es sich schliesslich nur noch um die Kongruenz\newpage
\SourceLeaf{20}\ArithmeticHead{13}{R}{INHALT DER ABSCHNITTE I—IV DER D. A.}
\noindent \(x^2\equiv a\ (\mathrm{mod.}\ p)\), wo \(a\) als durch \(p\) nicht teilbar gedacht werden darf. Jenachdem sie möglich ist oder nicht, heisst \(a\) nach \textsc{Euler}\footnote{Siehe Novi Comment. Acad. Petropol. \textbf{5} (1754/55) 1760, S. 13, Opera Omnia, Ser. I, vol. 2, S. 338.} quadratischer Rest oder Nichtrest (residuum oder nonresiduum) von \(p\), nach \textsc{Gauss} in Zeichen \(aRp\) oder \(aNp\). Nach \textsc{Legendre}\begingroup\renewcommand{\thefootnote}{**)}\footnote{Théorie des Nombres, 3\textsuperscript{me} éd. (1830) I, S. 197.}\endgroup\ bezeichnet man dieses verschiedene Verhalten, indem man dem Symbole \(\left(\frac{a}{p}\right)\) entsprechend den Wert \(+1\) oder \(-1\) beilegt. Es gibt der inkongruenten Zahlen \(a\) beider Arten gleichviel.

Bei weitem schwieriger aber erwies sich die Beantwortung der Frage in der zweiten der erwähnten Annahmen, der Frage: von welchen ungeraden Primzahlen \(p\) eine gegebene Zahl \(a\) quadratischer Rest bezw. quadratischer Nichtrest sei. Schon vor \textsc{Gauss} hatte \textsc{Euler}\begingroup\renewcommand{\thefootnote}{***)}\footnote{Opuscula Analytica I (1783), S. 135; Novi Comment. Acad. Petrop. \textbf{8} (1760/61) 1763, S. 105, Opera Omnia, Ser. I, vol. 2, S. 556.}\endgroup\ die Antwort auf diese Frage für die Werte \(a=-1\) und 3, \textsc{Lagrange}\begingroup\renewcommand{\thefootnote}{†)}\footnote{Nouv. Mém. de l’Acad. Berlin (1775) 1777, S. 351, 352, Oeuvres III, S. 791.}\endgroup\ für \(a=2,\,5,\,7\) gefunden, und \textsc{Gauss} selbst begründet zunächst seinerseits die so erhaltenen Sätze. Der berühmte \textsc{Euler}sche Satz, der nach \textsc{Gauss}’ oben erwähnter Aussage der Ausgangspunkt seiner tieferen arithmetischen Forschungen geworden ist, der Satz, dass \(-1\) quadratischer Rest von jeder Primzahl von der Form \(4n+1\), Nichtrest von jeder Primzahl von der Form \(4n+3\) ist, was sich einfach ausspricht in der Gleichung
\[
\left(\frac{-1}{p}\right)=(-1)^{\frac{p-1}{2}},
\]
folgt unmittelbar aus dem \textsc{Euler}schen Kriterium oder lässt sich mittels assoziierter Zahlen oder auch auf Grund des \textsc{Wilson}schen Satzes erhärten. Für den ebenso schönen Satz, dass 2 quadratischer Rest von jeder Primzahl von einer der Formen \(8n+1,\,8n+7\), Nichtrest von jeder Primzahl von einer der beiden Formen \(8n+3,\,8n+5\), dass also in \textsc{Legendre}scher Bezeichnung
\[
\left(\frac{2}{p}\right)=(-1)^{\frac{p^2-1}{8}}
\]
ist, gab \textsc{Gauss} an dieser Stelle seiner D. A. zwei zum Teil von einander abweichende einfache Begründungen. Nun aber zeigte sich, dass für den qua-\newpage
\SourceLeaf{21}\ArithmeticHead{14}{L}{BACHMANN, ÜBER GAUSS’ ZAHLENTHEORETISCHE ARBEITEN.}
\noindent dratischen Charakter einer ungeraden Primzahl \(q\) bezüglich einer andern solchen Primzahl \(p\) kein ähnlicher direkter Ausspruch angebbar, jener vielmehr nur umgekehrt durch den quadratischen Charakter von \(p\) in Bezug auf \(q\) bestimmbar ist. Das hier geltende eigentümliche Gesetz, das \textsc{Gauss} durch Induktion fand, war so auch vorher schon von \textsc{Legendre} ermittelt worden\footnote{Histoire de l’Acad. des Sciences, Paris 1785, S. 465.}, dem jedoch noch nicht gelang, es streng zu beweisen. Durch \textsc{Kronecker}\begingroup\renewcommand{\thefootnote}{**)}\footnote{Berliner Monatsberichte 1875, S. 267.}\endgroup\ aber ist zuerst bemerkt worden, dass \textsc{Euler}\begingroup\renewcommand{\thefootnote}{***)}\footnote{Opuscula Analytica I (1783), S. 64.}\endgroup\ als Entdecker des Gesetzes gelten darf, da er es in einer mit der von \textsc{Gauss} gegebenen wesentlich identischen Formulierung schon vor \textsc{Legendre} ausgesprochen hat. Diese \textsc{Gauss}sche Fassung besagt:

Wenn \(p\) eine Primzahl von der Form \(4n+1\), so ist \(+p\), wenn es eine Primzahl von der Form \(4n+3\) ist, so ist \(-p\) Rest oder Nichtrest von jeder Primzahl, welche positiv gedacht Rest resp. Nichtrest von \(p\) ist.

Mit Recht hat \textsc{Gauss} diesen Satz das theorema fundamentale der quadratischen Reste genannt, da es nicht nur in ihrer Theorie den eigentlichen Kern ausmacht, sondern auch für weitere Teile der höheren Arithmetik grundlegende Bedeutung hat. Kaum aber hat er schon voraussehen können, wie gross dessen Bedeutung für die ganze Entwicklung der Zahlentheorie werden, wie schon seine eigenen weiteren Beweise des Satzes, von denen wir zu sprechen haben werden, sodann die ungefähr 40 andern Beweise, welche wir späteren Forschern zu danken haben, mit immer neuen Gesichtspunkten auch neue Richtungen der Forschung herbeiführen, wie endlich das Bedürfnis, den analogen Satz für Potenzreste höheren Grades zu erledigen, dem schon \textsc{Gauss} selbst, später besonders \textsc{Kummer}\begingroup\renewcommand{\thefootnote}{†)}\footnote{Vergl. \textit{Über die allgemeinen Reziprozitätsgesetze u. s. w.} Abhandl. der Berliner Akademie 1859.}\endgroup\ und in unserer Zeit \textsc{Hilbert}\begingroup\renewcommand{\thefootnote}{††)}\footnote{Math. Annalen (1899) \textbf{51}, S. 1—127.}\endgroup\ erfolgreich ihre Bemühungen gewidmet, ganz neue Welten mathematischer Ideen uns eröffnen und gewinnen lassen würde!
